\documentclass[10pt,conference]{IEEEtran}

\usepackage{graphicx}
\usepackage{amsmath,amssymb}
\usepackage{booktabs}
\usepackage{array}
\usepackage{xcolor}
\usepackage{hyperref}
\usepackage{url}
\usepackage{enumitem}
\usepackage{placeins}

\graphicspath{{figures/}}

\hypersetup{
    colorlinks=true,
    linkcolor=blue,
    citecolor=blue,
    urlcolor=blue
}

% ── Helper macros ─────────────────────────────────────────────────────────────

\newcommand{\figplaceholder}[1]{%
\fbox{%
\begin{minipage}[c][1.8in][c]{0.9\linewidth}
\centering
Placeholder / missing figure:\\[0.5em]
\texttt{\detokenize{#1}}
\end{minipage}}}

\newcommand{\reportfigure}[4]{%
\begin{figure*}[t]
\centering
\IfFileExists{figures/#1}{%
    \includegraphics[width=#2]{figures/#1}
}{%
    \figplaceholder{figures/#1}
}
\caption{#3}
\label{#4}
\end{figure*}
}

\newcommand{\reporttable}[3]{%
\begin{table*}[t]
\centering
\IfFileExists{figures/tables_tex/#1}{%
    \resizebox{\textwidth}{!}{\input{figures/tables_tex/#1}}
}{%
    \fbox{%
    \begin{minipage}[c][1.4in][c]{0.9\linewidth}
    \centering
    Placeholder / missing table:\\[0.5em]
    \texttt{\detokenize{figures/tables_tex/#1}}
    \end{minipage}}
}
\caption{#2}
\label{#3}
\end{table*}
}

% ── Title ─────────────────────────────────────────────────────────────────────

\title{Component-Hierarchy and Merge-Tree Proxy Analysis\\
for Wind-Field Super-Resolution}

\author{
\IEEEauthorblockN{Arjun Dadhwal}
\IEEEauthorblockA{
Department of Computer Science\\
Tulane University\\
}
}

\begin{document}

\maketitle

\begin{abstract}
This report analyzes scalar wind-speed super-resolution outputs from bicubic
interpolation, a CNN model, and a GAN model. The goal is to understand whether
topology-based evaluators, especially merge trees, capture structural differences
that are not explained by conventional pixelwise, distributional, or physics-domain
metrics alone. Prior inspection showed that PSNR and SSIM consistently favor CNN
outputs, while persistence diagrams strongly favor GAN outputs and merge trees
select GAN in only a small subset of samples. To investigate this discrepancy, we
compute superlevel-set component counts, area-filtered component statistics,
exceedance curves, and branch-lifetime proxy statistics across the 168-sample
evaluation set. The results suggest that GAN outputs often recover additional
component complexity, but can also over-fragment the scalar field. CNN outputs are
smoother and often closer in direct or aggregate physical error, but may suppress
high-frequency structure. Importantly, adjacent-control cases show that merge-tree
decisions are not fully explained by component counts or simple branch summaries,
motivating future work on topology-aware losses for scientific super-resolution.
\end{abstract}

% ── 1. Introduction ───────────────────────────────────────────────────────────

\section{Introduction}

Scientific super-resolution models are often evaluated using direct error metrics
such as PSNR, SSIM, MAE, and RMSE. These metrics are useful for measuring pointwise
fidelity, but they may not fully capture whether reconstructed scalar fields preserve
physically and topologically meaningful structure. This issue is especially important
for wind-field super-resolution, where small-scale ridges, high-speed regions, and
connected structures can matter for downstream interpretation.

This project extends the evaluation of PhIRE-style wind-field super-resolution
\cite{stengel2020adversarial} by comparing CNN, GAN, and bicubic outputs using both
traditional metrics and topology-inspired descriptors. The initial motivation was the
perception-distortion tradeoff \cite{blau2018perception}: GAN outputs may appear
sharper and more realistic, while CNN outputs may be more pointwise accurate. However,
for scientific fields, sharper structure is not automatically better. A sharper feature
may correspond to a physically meaningful ridge, or it may be a hallucinated artifact.

Earlier analysis showed that SSIM selected CNN for all 168 samples, while bottleneck
persistence-diagram distance selected GAN for nearly all samples. Merge-tree distance
was more selective, selecting GAN in 20 samples and CNN in the remaining 148. This
report focuses on understanding what is special about those 20 merge-tree-GAN cases,
whether simple component-count statistics explain them, and whether richer
hierarchy-aware proxies provide better insight.

% ── 2. Data and Compared Methods ─────────────────────────────────────────────

\section{Data and Compared Methods}

The evaluation set consists of 168 paired low-resolution and high-resolution
wind-field samples. Each sample contains vector wind components $(u,v)$. All topology
and component analyses in this report are computed on scalar wind-speed magnitude,
\begin{equation}
s(x,y) = \sqrt{u(x,y)^2 + v(x,y)^2}.
\end{equation}

We compare four fields per sample:
\begin{enumerate}[leftmargin=*]
    \item Ground truth high-resolution speed field, denoted GT.
    \item Bicubic interpolation baseline.
    \item CNN super-resolution output.
    \item GAN super-resolution output.
\end{enumerate}

The bicubic baseline is included as a non-learned interpolation control. This is
useful because it shows how much structure is already recoverable by smooth
interpolation and how much additional behavior is introduced by learned super-resolution.

\reportfigure
{fig01_pipeline_placeholder.pdf}
{0.95\textwidth}
{Pipeline overview. The evaluation begins with low-resolution wind inputs and
high-resolution ground truth fields. Bicubic, CNN, and GAN outputs are converted to
scalar speed magnitude. The resulting scalar fields are evaluated using direct error
metrics, physics/domain metrics, persistence diagrams, merge trees, component counts,
and merge-tree proxy statistics. This is a placeholder figure to be drawn manually.}
{fig:pipeline}

% ── 3. Evaluation Metrics ────────────────────────────────────────────────────

\section{Evaluation Metrics}

The analysis combines several metric families.

\subsection{Direct Metrics}
Direct metrics compare reconstructed fields to GT pointwise. These include PSNR,
SSIM, MAE, and RMSE. They generally reward smooth, conservative reconstructions and
penalize high-frequency differences.

\subsection{Physics and Domain Metrics}
Physics/domain metrics summarize physically relevant scalar-field behavior. These
include wind-power-density errors, gradient statistics, power spectral density
statistics, Wasserstein distances, and exceedance/tail behavior. These metrics help
distinguish whether a model preserves meaningful physical distributions, gradients,
and high-speed events.

\subsection{Topological Metrics}
Two topology-based validators are considered:
\begin{enumerate}[leftmargin=*]
    \item Persistence diagrams, compared using bottleneck distance.
    \item Merge trees, compared using a merge-tree distance.
\end{enumerate}

Persistence diagrams summarize the birth and death of topological features but
discard much of the spatial and hierarchical organization. Merge trees retain more
information about how superlevel-set components appear and merge as the threshold
changes. This makes merge trees a useful candidate for identifying whether
GAN-generated features are structurally meaningful or merely sharp artifacts.

% ── 4. Selector-Level Summary ────────────────────────────────────────────────

\section{Selector-Level Summary}

Table~\ref{tab:selector_counts} summarizes the high-level winner counts over the
168-sample evaluation set.

\begin{table}[t]
\centering
\caption{All-sample winner counts across the 168-sample evaluation set.}
\label{tab:selector_counts}
\begin{tabular}{lccc}
\toprule
Selector / validator & CNN wins & GAN wins & Ties / unavailable \\
\midrule
PSNR$_{uv}$             & 168 & 0   & 0 \\
SSIM                    & 168 & 0   & 0 \\
Merge tree (MT)         & 148 & 20  & 0 \\
Bottleneck PD           & 2   & 166 & 0 \\
Physics group majority  & 168 & 0   & 0 \\
\bottomrule
\end{tabular}
\end{table}

This result motivates the rest of the report. Direct metrics and SSIM strongly favor
CNN, while persistence diagrams overwhelmingly favor GAN. Merge trees are more
selective: they usually favor CNN but select GAN in 20 cases. This suggests that
merge trees may be capturing a different structural signal than either direct error
metrics or persistence diagrams.

\reportfigure
{fig02_three_method_metric_summary.pdf}
{0.95\textwidth}
{Three-method metric winner summary for bicubic, CNN, and GAN. This figure compares
which method wins across the metric suite. It provides a baseline view of how direct,
distributional, tail, and physics-domain metrics distribute across the three
reconstruction methods.}
{fig:three_method_summary}

% ── 5. Superlevel-Set Component Counts ───────────────────────────────────────

\section{Superlevel-Set Component Counts}

A simple way to probe topology is to count connected components in superlevel sets,
\begin{equation}
S_{\tau} = \{(x,y) : s(x,y) \geq \tau\}.
\end{equation}
For each threshold $\tau$, the number of connected components provides a coarse
description of scalar-field structure. However, raw component counts do not capture
how components are nested or how they merge across thresholds.

The component-count analysis uses both fixed thresholds and percentile thresholds:
\begin{equation}
\tau \in \{5,10,15,p90,p95,p99\}.
\end{equation}

\reporttable
{table_component_count_summary.tex}
{Summary of superlevel-set component-count behavior across the 168-sample evaluation
set. This table is generated from \texttt{component\_count\_summary.csv}.}
{tab:component_count_summary}

\reportfigure
{fig03_component_count_threshold_heatmap.pdf}
{0.75\textwidth}
{Component-count winner counts by threshold. GAN tends to better match GT component
counts at lower thresholds, while CNN often wins at higher or percentile-based
thresholds.}
{fig:component_heatmap}

The component-count results reveal a useful but incomplete pattern. GAN frequently
matches GT component counts at low thresholds, suggesting that it restores broad
component complexity that CNN and bicubic interpolation smooth away. At higher
thresholds, CNN often becomes closer to GT, suggesting that GAN may over-fragment
high-speed regions or introduce excessive small-scale structure.

% ── 6. Special Component-Count Cases ─────────────────────────────────────────

\section{Special Component-Count Cases}

The most important subset consists of the 20 samples where merge trees favor GAN.
These cases test whether GAN is doing something structurally meaningful, or whether
the merge-tree distance is being influenced by sharp but potentially misaligned texture.

\reportfigure
{fig04_mt_gan_signature_matrix.pdf}
{0.8\textwidth}
{Component-count winner signatures for the 20 MT-GAN cases. Each row is one MT-GAN
sample, and each column is a threshold. The matrix shows whether bicubic, CNN, or GAN
has the closest component count to GT at each threshold.}
{fig:mt_gan_signature_matrix}

The signature matrix shows that MT-GAN samples are not uniform. Some samples show
strong GAN support across thresholds, while others only support GAN at low thresholds
or show mixed behavior. Therefore, MT selecting GAN cannot be reduced to a single
rule such as ``GAN has more components.''

% ── 7. Adjacent Control Clusters ─────────────────────────────────────────────

\section{Adjacent Control Clusters}

Adjacent sample IDs are useful controls because they may correspond to nearby
timesteps and visually similar wind-field regimes. If MT selects GAN for one sample
but CNN for its neighbors, then the cause may be a subtle local or hierarchical
difference rather than a large global change in field complexity.

We inspect three adjacent clusters:
\begin{itemize}[leftmargin=*]
    \item Samples 10--13, where sample 12 is MT-GAN.
    \item Samples 76--78, where sample 77 is MT-GAN.
    \item Samples 90--93, where sample 92 is MT-GAN while nearby samples remain MT-CNN.
\end{itemize}

\reportfigure
{fig05_adjacent_cluster_10_13_component_curves.pdf}
{0.95\textwidth}
{Adjacent cluster 10--13. Sample 12 is an MT-GAN case, while neighboring samples
provide local controls.}
{fig:cluster_10_13}

\reportfigure
{fig05_adjacent_cluster_76_78_component_curves.pdf}
{0.95\textwidth}
{Adjacent cluster 76--78. Sample 77 is an MT-GAN case, surrounded by neighboring
control samples.}
{fig:cluster_76_78}

\reportfigure
{fig05_adjacent_cluster_90_93_component_curves.pdf}
{0.95\textwidth}
{Adjacent cluster 90--93. Sample 92 is an MT-GAN case, while nearby samples are
useful controls because several also have GAN-favoring domain metrics.}
{fig:cluster_90_93}

These adjacent clusters show that component-count curves alone do not fully explain
MT decisions. Neighboring samples can have similar component-count behavior while MT
selects different models. This suggests that MT is sensitive to information discarded
by marginal component counts, such as spatial arrangement, branch nesting, or merge
order.

% ── 8. Rare Topology-CNN Controls ────────────────────────────────────────────

\section{Rare Topology-CNN Controls}

Samples 162 and 163 are rare cases where both bottleneck persistence diagrams and
merge trees favor CNN. These are important negative controls because they show that
topological descriptors do not automatically reward GAN sharpness.

\reportfigure
{fig08_rare_topology_cnn_controls_161_164.pdf}
{0.95\textwidth}
{Rare topology-CNN control neighborhood. Samples 162 and 163 are rare cases where
both PD and MT favor CNN. Samples 161 and 164 provide adjacent context.}
{fig:rare_topology_cnn}

These controls help distinguish between two possibilities. If topology always preferred
GAN, then topology would simply be rewarding sharper or noisier fields. Instead, the
rare topology-CNN samples show that topological validators can reject GAN structure
when it is not aligned with GT.

% ── 9. Dense Component-Hierarchy Proxies ─────────────────────────────────────

\section{Dense Component-Hierarchy Proxies}

The six-threshold component-count analysis is informative but coarse. To better
approximate merge-tree behavior, we compute dense threshold curves using 40
GT-percentile thresholds from p1 to p99. We also repeat component counts after
filtering out small components using minimum area thresholds:
\begin{equation}
a_{\min} \in \{1,5,10,25,50,100\}.
\end{equation}

This analysis helps separate meaningful components from small speckle-like fragments.
If GAN only wins when $a_{\min}=1$, the improvement may be due to small noisy
components. If GAN remains competitive at larger area filters, then its additional
structure is more likely to be spatially meaningful.

\reportfigure
{fig09_dense_component_curves_representative_samples.pdf}
{0.95\textwidth}
{Dense area-filtered component-count curves for representative samples. These curves
use 40 GT-percentile thresholds and a minimum component area of 10 pixels.}
{fig:dense_component_curves}

% ── 10. Branch-Lifetime Merge-Tree Proxies ───────────────────────────────────

\section{Branch-Lifetime Merge-Tree Proxies}

Component counts still do not encode the hierarchy of merges. Therefore, we compute
branch-lifetime proxy statistics using a descending-threshold sweep. As the threshold
decreases, new superlevel components are born near local maxima, grow, and merge into
older components. When a younger component merges into an older component, we record
a branch lifetime:
\begin{equation}
\mathrm{persistence} = s_{\mathrm{birth}} - s_{\mathrm{death}}.
\end{equation}

This is not an exact merge-tree distance, but it approximates several
hierarchy-aware quantities:
\begin{itemize}[leftmargin=*]
    \item total branch count,
    \item number of branches above persistence thresholds,
    \item mean, median, and maximum persistence,
    \item total persistence,
    \item top-5 persistence sum,
    \item persistence entropy.
\end{itemize}

\reporttable
{table_merge_tree_proxy_summary.tex}
{Method-level summary of merge-tree proxy errors. This table is generated from
\texttt{merge\_tree\_proxy\_summary.csv}. Lower values indicate closer agreement with
GT for the displayed proxy.}
{tab:merge_tree_proxy_summary}

\reportfigure
{fig10_hierarchy_proxy_summary.pdf}
{0.95\textwidth}
{Hierarchy-proxy summary by case family. The top panel compares mean branch-count
ratios relative to GT. The bottom panel shows GAN mean-persistence bias relative
to GT.}
{fig:hierarchy_proxy_summary}

The branch-lifetime proxies suggest a more nuanced interpretation. GAN often restores
component complexity, but it can also produce too many branches relative to GT. CNN
tends to smooth structure and may under-produce branches, but it can still be closer
to GT in branch-count error or aggregate persistence. This supports the idea that the
CNN/GAN tradeoff is not simply pointwise accuracy versus visual sharpness; it also
involves structural over-fragmentation versus structural smoothing.

% ── 11. Interpretation ───────────────────────────────────────────────────────

\section{Interpretation}

The main finding is that merge-tree decisions cannot be fully explained by simple
component counts or global branch summaries. The component-count analysis shows that
GAN is often closer to GT at low thresholds, while CNN is often closer at high or
percentile thresholds. The branch-proxy analysis suggests that GAN may over-fragment
the field, producing many additional branch events. However, adjacent control clusters
show that MT-GAN selections are not isolated by these marginal summaries alone.

This leads to the following interpretation:
\begin{quote}
GAN outputs often recover additional topological complexity that CNN suppresses, but
this added complexity is not always hierarchically or spatially aligned with GT.
Merge trees appear to be sensitive to more than the number of components or branches;
they may respond to where structures occur, how they merge, and how persistent
structures are nested.
\end{quote}

This interpretation is important because it suggests why persistence diagrams and
merge trees can disagree. Persistence diagrams summarize birth-death pairs while
discarding much of the explicit spatial and hierarchical organization. Merge trees
preserve more of this organization, making them more selective when evaluating
GAN-generated structures.

% ── 12. Implications for Topology-Aware Training ─────────────────────────────

\section{Implications for Topology-Aware Training}

The analysis motivates a topology-aware scientific super-resolution model. A useful
training objective should not simply encourage sharper output. Instead, it should
preserve meaningful multiscale structure while avoiding GAN-style over-fragmentation.

A possible future objective is:
\begin{equation}
\mathcal{L}
=
\lambda_{\mathrm{pix}}\mathcal{L}_{\mathrm{pix}}
+
\lambda_{\mathrm{phys}}\mathcal{L}_{\mathrm{phys}}
+
\lambda_{\mathrm{topo}}\mathcal{L}_{\mathrm{topo}},
\end{equation}
where $\mathcal{L}_{\mathrm{pix}}$ preserves pointwise accuracy,
$\mathcal{L}_{\mathrm{phys}}$ preserves physical/domain statistics, and
$\mathcal{L}_{\mathrm{topo}}$ encourages structure-aware agreement with GT.

Exact merge-tree distances may be difficult to differentiate directly. Therefore,
the first practical step may be to use differentiable topology-inspired proxies, such
as soft exceedance losses, soft component-area losses, persistence-inspired penalties,
or learned merge-tree similarity models such as MTNN \cite{qin2024mtnn}. Exact PD
and MT distances can remain post hoc evaluators.

% ── 13. Limitations ──────────────────────────────────────────────────────────

\section{Limitations}

This report uses branch-lifetime and component-hierarchy proxies rather than exact
merge-tree visualization or exact differentiable merge-tree losses. These proxies
are useful for diagnosis, but they compress the full merge-tree structure. They do
not fully encode branch matching, spatial embedding, or detailed merge order.

The visual interpretation is also challenging. Wind-speed scalar fields contain many
small-scale structures, and GAN outputs can appear visually plausible even when they
are topologically over-fragmented. Therefore, the strongest evidence comes from
combining multiple views: direct error maps, component-count curves, branch proxies,
adjacent controls, and rare topology-CNN controls.

% ── 14. Conclusion ───────────────────────────────────────────────────────────

\section{Conclusion}

This report extends the initial CNN/GAN comparison by adding bicubic interpolation,
component-count analysis, adjacent controls, and merge-tree proxy statistics. The
results show that GAN often restores component complexity but may also over-fragment
the scalar field. CNN is more conservative and often better under direct and aggregate
physical metrics, but it may suppress meaningful high-frequency structure. Merge trees
are more selective than persistence diagrams and do not appear to follow component
counts alone.

The most important outcome is that topology-aware evaluation reveals structural
behavior that conventional metrics do not fully capture. This motivates the next phase
of the project: developing a topology-aware super-resolution model that balances
CNN-like pointwise accuracy with better preservation of meaningful multiscale structure.

\bibliographystyle{IEEEtran}
\bibliography{main}

\end{document}